\section{Discussion and Limitations}
\label{sec:discussion}
\subsection{What the experiment establishes}
The main empirical result is a net performance shortfall relative to SPY buy-and-hold. Three measured components explain it. First, squared-error prediction is slightly worse than the historical-mean reference, and direction accuracy does not exceed the unconditional positive-return rate. Second, the strategy omits intervals with positive average equity returns, lowering gross growth. Third, 246 one-way trade units create a substantial fee burden. These components are consistent with the observed combination of lower volatility, lower CAGR, and lower net Sharpe.

The financial and statistical comparisons answer different questions. The negative $R^2_{\mathrm{OOS}}$ concerns squared error; the portfolio result concerns weighted returns and costs. Here both offer little support for the selected rule, but the connection is empirical rather than an identity. The strategy also selects an invested subset with a higher mean asset return than the excluded subset. That association is insufficient to offset foregone positive exposure and implementation costs. The joint evidence is more informative than either a hit rate or a drawdown ranking alone.

Ridge has a clear role within the workflow: it constrains coefficient magnitudes in a correlated-input regression. The rolling window and monthly schedule determine how much history informs each decision. Validation then chooses among four implementations of that estimator. This separation keeps the interpretation precise. The experiment evaluates a regularized trading procedure, not a theorem that shrinkage or more features must produce an investment advantage.

\subsection{Limitations}
\label{sec:limitations}
The strategy was designed in November 2026 using already historical market data. The implemented configuration and selection were fixed before its initial download, but the study is not an externally preregistered prospective experiment. Selecting the largest validation Sharpe can favor noise even with four candidates. The fall from validation Sharpe 1.206 to test Sharpe 0.481 can reflect selection optimism, changing relationships, or sampling variation; the present design does not distinguish those explanations.

The final evaluation concerns one ETF and four years. Sharpe ratios and annual results are point estimates: no dependence-robust confidence intervals, formal superiority tests, or factor-alpha estimates are included. The gross Sharpe difference of approximately 0.007 is consequently treated as numerical similarity, not evidence of outperformance. Conditional-return and coefficient summaries are descriptive associations among correlated observations and inputs.

Execution assumes the observed next opening price and a fixed proportional fee. Opening-auction depth, market impact, partial fills, and actual brokerage financing are not modeled. Cash earns zero, so the reported Sharpe is a zero-rate measure rather than realized excess performance over Treasury bills. Dividend receivables are credited and reinvested at the ex-date exit open instead of their payment date. The supported history has no splits. Vendor revisions can alter a later download, and exact numerical reproduction depends on the recorded input snapshot.

The constant-exposure diagnostic uses the observed test-average allocation and ignores fractional rebalancing costs. It is an ex post attribution calculation, not an investable benchmark. A risk-matched strategy would require its own advance allocation rule and accounting. Likewise, no-trade regions, continuous weights, state interactions, other training windows, and additional assets are outside the evaluated strategy. These are changes to the research design rather than results of this study.

\section{Conclusion}
\label{sec:conclusion}
The selected Ridge rule converts five trailing price features into daily SPY or cash positions under explicit timing, training, and cost conventions. In 2022--2025 it earns \RidgeCAGR{} net CAGR and a \RidgeSharpe{} Sharpe ratio, below buy-and-hold's \BHCAGR{} and \BHSharpe{}. Its smaller full-period drawdown accompanies lower equity participation. Weak forecast performance, positive returns during excluded intervals, and frequent full-position switches jointly explain the economic result. The project demonstrates the complete path from a plausible signal specification to a validated, reproducible trading evaluation, while preserving the finding that this implementation does not improve net performance.

\section{Reproducibility and Contributions}
\label{sec:repro}
The research implementation, configuration, and published result files are available in the accompanying repository \citep{implementation}. Training and accounting are implemented in \code{project2/core.py}; data checks in \code{project2/data.py}; validation and final evaluation in \code{project2/pipeline.py}. The saved CSV and JSON outputs produce the report's tables and figures. Appendix~\ref{app:reproduction} provides platform-specific environment activation, platform-independent analysis commands, package versions, and the artifact map.

The recorded clean-environment rerun used the same input snapshot in a fresh working directory and reproduced the validation grid, test forecasts, fit audits, and headline metrics byte for byte. Eleven integrity checks cover label maturity, future-data perturbations, scaling, execution and distribution alignment, costs, data validity, and drawdown. Full fit records and individual test cases remain in the code package; the report retains representative records and the complete performance comparisons.

\subsection*{Author contributions}
\noindent\begin{tabularx}{\textwidth}{@{}p{.27\textwidth}X@{}}
\toprule
Member & Contribution \\
\midrule
\MemberOneName & \MemberOneContribution \\
\MemberTwoName & \MemberTwoContribution \\
\MemberThreeName & \MemberThreeContribution \\
\MemberFourName & \MemberFourContribution \\
\bottomrule
\end{tabularx}

\subsection*{AI assistance}
GPT-6.1 Sol (OpenAI) was used for code review and debugging, and for report organization and language editing. The authors verified the reported results and take responsibility for the analysis and submitted work.
